% --------------------- DATA & METHODOLOGY ----------------------------
\section{Data and Methodology}
\label{sec:data}

\subsection{Data and forecasting design}

The investable universe is OSEAX. For each firm and each semi-annual rebalancing event from January 2006 to March 2023, I construct snapshots 30, 60 and 100 trading days before the effective date. The predictors are intentionally restricted to variables that an external investor can observe before the announcement: trading status, 12-month electronic-order-book turnover trimmed for extreme days, free-float-adjusted market capitalisation, ICB super-sector, and four lags of OSEBX membership. This restriction is central to the design. The model is not trying to reverse-engineer hidden discretionary decisions; it is trying to process the same public rule-book inputs faster and more consistently than a manual screen.

Price returns use VWAP series from Bloomberg and Refinitiv DataStream, which better approximate institutional execution than closing prices \parencite{madhavan-vwap,bloomberg-terminal,lseg-refinitiv}. Norwegian Fama--French factors and the risk-free rate are from \textcite{bernt-fama-french} and \textcite{bernt-risk-free}. After cleaning, the three horizon panels contain 5{,}983, 5{,}924 and 5{,}891 firm-event observations. All reported performance is out of sample: for each test rebalancing, the training set contains only events whose cut-off dates are strictly earlier.

\subsection{Models and conditional threshold}

The first classifier is a binomial logistic GLM,
\begin{equation}
\Pr(Y_{i,t}=1 \mid X_{i,t}) = \frac{\exp(\beta^\top X_{i,t})}{1+\exp(\beta^\top X_{i,t})},
\label{eq:glm}
\end{equation}
where $Y_{i,t}=1$ means that firm $i$ is in OSEBX after rebalancing $t$. The second is XGBoost, a regularised boosted-tree model that can capture non-linear interactions between size, liquidity, sector and past membership \parencite{xgboost-a-scaleable-tree-boosting-system}. Hyperparameters are fixed across folds to avoid turning the backtest into a tuning exercise.

The main methodological issue is class persistence. Most current constituents remain in the index and most non-constituents remain out, so ordinary 0--1 loss rewards a model that predicts no additions or deletions. I therefore replace the standard threshold $\theta$ with a threshold conditioned on prior membership:
\begin{equation}
\theta^{\star}_{i,t} =
\begin{cases}
\theta + \varepsilon, & Y_{i,t-1}=1,\\
\theta - \varepsilon, & Y_{i,t-1}=0,
\end{cases}
\qquad
\hat{Y}_{i,t}=\mathbf{1}\{\hat{p}_{i,t}>\theta^{\star}_{i,t}\}.
\label{eq:cppt}
\end{equation}
With $\varepsilon>0$, a current constituent needs stronger evidence to be classified as staying, while a current non-constituent needs less evidence to be classified as entering. The mechanism is deliberately simple: it trades a small amount of calibration accuracy for more addition and deletion signals. I set $\theta=0.6$ and $\varepsilon=0.2$ on the training data. In the results, XGBB denotes standard XGBoost and XGBC denotes XGBoost with this conditional threshold.

\subsection{Portfolio tests}

For each rebalancing event, predicted additions $\mathcal{A}_t$ are bought and predicted deletions $\mathcal{D}_t$ are shorted at the forecast date, then closed at the effective date. Trade returns are equally weighted:
\begin{equation}
r_t^A =
\frac{1}{|\mathcal{A}_t \cup \mathcal{D}_t|}
\sum_{i \in \mathcal{A}_t \cup \mathcal{D}_t} r_i,
\qquad
r_i^{long}=\frac{p_1-p_0}{p_0},\quad
r_i^{short}=\frac{p_0-p_1}{p_0}.
\label{eq:retdef}
\end{equation}
The active simulation applies 4 bps per leg in transaction costs \parencite{dnb-transaction-costs}, a 4\% annual short-borrow charge \parencite{nordnet-short-interest-rate}, and a forced buy-back cap at a 100\% short loss. Outside active event windows the capital is invested in OSEBX.

The enhanced-index version combines OSEBX with the active sleeve:
\begin{equation}
r_t^E=s r_t^A+(1-s)r_t^{OSEBX},
\label{eq:enhanced}
\end{equation}
where $s$ is the largest realised weight that keeps annual tracking error below 2\%, a common enhanced-index risk budget \parencite{ssg-tracking-error-enhanced}. In this version, shorted deletion names are treated as sales from the passive inventory, so the short-borrow charge is removed. Risk-adjusted performance is measured with CAPM and Norwegian Fama--French three-factor regressions over January 2010 to May 2022, the longest period with factor data.
